\documentclass[10pt,a4paper]{amsart}
\usepackage[margin=19mm]{geometry}
\usepackage{amsmath,amssymb,mathtools}
\usepackage[scr=rsfs,cal=euler]{mathalfa}
\usepackage{xfrac}
\usepackage[unicode,hidelinks,bookmarks=false]{hyperref}
\newcommand{\ie}{i.e.\ }
\newcommand{\cf}{cf.\ }
\newcommand{\resp}{respectively\ }
\newcommand{\Subsection}{\subsection}
\providecommand{\nobreakdash}{\nobreak\hbox{-}}
\setlength{\emergencystretch}{2em}
\multlinegap0pt
\usepackage{bm}
\usepackage{bbm}
\usepackage{dsfont}
\usepackage{xspace}
\usepackage{cancel}
\usepackage{mathtools}
\usepackage{amsthm}
\makeatletter
\@ifundefined{theorem}{\newtheorem{theorem}{Theorem}}{}
\@ifundefined{lemma}{\newtheorem{lemma}[theorem]{Lemma}}{}
\makeatother
\makeatletter
\@ifundefined{claimproof}{\newenvironment{claimproof}
  {\renewcommand{\qedsymbol}{\ensuremath{\blacksquare}}\proof}}{}
\makeatother
\title[The Tur\'{a}n number of the Cartesian product of a star and ]{The Tur\'{a}n number of the Cartesian product of a star and an edge}
\author{Xiamiao Zhao, Xin Cheng, Cheng Chi, Ervin Gy\H{o}ri, Casey Tompkins, Yichen Wang}
\date{Source: arXiv:2604.11366v1 (CC BY 4.0)}
\begin{document}
\maketitle

\noindent\emph{Source and licence.} The Tur\'{a}n number of the Cartesian product of a star and an edge. Version arXiv:2604.11366v1 (CC BY 4.0), distributed under the Creative Commons Attribution 4.0 International licence, as recorded in the arXiv licence metadata for this exact version and shown on the abstract page https://arxiv.org/abs/2604.11366v1. Full rights notice: licenses/arxiv-2604.11366-CC-BY-4.0.txt.\par\smallskip

\noindent\textit{Editorial scope.} Counted unit: the Lemma whose source label is \texttt{lem: lower bound of number of C4 for Bt on bip} in the article (source0.tex lines 884--891, source label \texttt{lem: lower bound of number of C4 for Bt on bip}), delivered together with the article's own complete original proof of it (source0.tex lines 893--996, one \texttt{proof} environment of 104 lines). No mathematics is written, repaired or re-derived: statement and proof are the article's own text line for line. Required context retained uncounted: lem: lower bound of number of C4 for Bt (lines 740--748). Every definition, display and label of those lines is retained unchanged. Omitted from this package and not redistributed: 1--739, 749--883, 892--892, 997--1938. Nothing inside a retained excerpt is omitted or altered: every retained body is delivered exactly as the article's lines are written, so no omission receipt of any kind accompanies this package. Citations: the retained excerpts carry no citation command at all, so no bibliography entry of the article is transcribed and no reference is redistributed. Reference adaptation: none was needed and none was made; every reference inside the delivered unit resolves inside it, so no unresolved reference or citation is produced. Printed slips kept literal: no printed word, symbol, hypothesis or number of the counted statement or of its proof was altered. Layout and class: the article's own class is \texttt{article} with packages including amsthm, amsmath, amssymb, tikz, xcolor, graphicx, mathtools, enumitem, booktabs, float, url, aligned-overset, fontenc, caption; the wrapper is the lane's pinned \texttt{amsart} editorial wrapper at 10pt on A4, which restates the theorem-like environments the retained text needs and adds the article's own macro definitions that the retained text needs (none), with extra wrapper packages (bm, bbm, dsfont, xspace, cancel, mathtools). The delivered numbering is the unit's own. Assets: no retained span contains a figure environment, a graphics-inclusion command, a PDF-page inclusion, a table or a TikZ picture; no image file is copied into this package, the package declares no asset dependency and no image file is redistributed with it. Licence: version arXiv:2604.11366v1 of this article is distributed under the Creative Commons Attribution 4.0 International licence, as recorded in the arXiv licence metadata for this exact version and shown on the abstract page https://arxiv.org/abs/2604.11366v1. A full rights notice is delivered at licenses/arxiv-2604.11366-CC-BY-4.0.txt. Characters: every retained excerpt is pure ASCII, so the delivered engine, XeLaTeX with its default Latin Modern text font, reports no missing character.
        \begin{lemma}\label{lem: lower bound of number of C4 for Bt}
            Let $t$ be a positive integer and $G$ be a $B_t$-free graph on $n$ vertices. When $n$ is sufficiently large, we have
            % \begin{equation*}
            %     N(C_4, G) \geq \frac{1}{4} \binom{t}{2} (\sum_{u \in V} d(u)^2 - 2(4t-1) e(G) - (t-1) n^2).
            % \end{equation*}
            \begin{equation*}
                \sum_{\substack{S \subseteq V(G), |S| = t \\ S \text{ is good}}} d(S)(d(S) - 1) \geq \sum_{u \in V(G)} d(u)^2 - 2(4t-1) e(G) - (t-1) n^2.
            \end{equation*}
        \end{lemma}

        \begin{lemma}\label{lem: lower bound of number of C4 for Bt on bip}
            Let $t$ be a positive integer and $G$ be a $B_t$-free bipartite graph with parts $V_1$ and $V_2$, and $|V_i| = n_i$.
            When $n$ is sufficiently large, we have
            \begin{equation*}
                N_i(C_4, G) \geq \frac{1}{2} \binom{t}{2} \left(\sum_{u \in V_i} d(u)^2 - (2t-1) e(G) - (t-1) n_j^2 \right),
            \end{equation*}
            where $i \neq j$, $i, j \in \{1, 2\}$.
        \end{lemma}

        \begin{proof}
            For a vertex $v \in V_1$, let $\mathcal{H}_v$ be the auxiliary hypergraph defined in the proof of Lemma \ref{lem: lower bound of number of C4 for Bt}.
            % The vertex set of $\mathcal{H}_v$ is $V(G) \setminus \{v\}$, and the hyperedge set is $\{h: h = N(v) \cap N(w) \cup \{w\}, w \in V_1 \setminus \{v\}\}$.
            % By the definition of $\mathcal{H}_v$, each hyperedge of size $s + 1$ corresponds to $s$ edges which incident to some vertices of $N(v)$.
            Since $G$ is a bipartite graph, we have
            \begin{equation}\label{eq: sum of dv bip}
                \sum_{u \in N(v)} d(u) = d(v) + \sum_{h \in E(\mathcal{H}_v)} (|h| - 1).
            \end{equation}
            % where the first item on the right is the number of edges between $v$ and $N(v)$, the second item is the number of edges between $N(v)$ and $V_1$.

            Let $\mathcal{H}_{v, i}$ ($1 \leq i \leq 3$) be the subhypergraph of $\mathcal{H}_v$ as defined in the proof of Lemma \ref{lem: lower bound of number of C4 for Bt}.
            % Let $\mathcal{H}_{v, 1}$ denote the subhypergraph of $\mathcal{H}_v$ which consisting of hyperedges with size at most $t$, $\mathcal{H}_{v, 2}$ denote the subhypergraph of $\mathcal{H}_v$ which consisting of hyperedges with size $t+1$, and $\mathcal{H}_{v, 3}$ denote the subhypergraph of $\mathcal{H}_v$ which consisting of hyperedges with size at least $t+2$.
            % Thus $\sum_{h \in E(\mathcal{H}_v)} (|h| - 1) = \sum_{i = 1}^{3} \sum_{h \in E(\mathcal{H}_{v, i})} (|h| - 1)$.

            % \begin{claimthm}
            %     For each vertex $u \in N(v)$, it can be contained in at most $t-1$ hyperedges of $\mathcal{H}_{v, 3}$.
            % \end{claimthm}
            % \begin{proof}
            %     Suppose not, there are at least $t$ hyperedges of $\mathcal{H}_{v, 3}$, namely $h_1, h_2, \dots, h_t$, contain $u$.
            %     Let $w_i$ be the vertex in $h_i$ but not in $N(v)$.
            %     Since the size of $h_i$ is at least $t+2$, we can choose vertices $u_i \neq w_i$ from $h_i$ greedily such that $u_i \neq u_j$ for $1 \leq i, j \leq t$.
            %     Then there are $t$ copies of $C_4$, namely $v u w_1 u_1$, $v u w_2 u_2$, \dots, $v u w_t u_t$, which share a common edge $vu$ and form a $B_t$, a contradiction.
            % \end{proof}
            As in the proof of Lemma \ref{lem: lower bound of number of C4 for Bt}, we have 
            \begin{equation}\label{eq: lem2.2 1}
                \sum_{h \in E(\mathcal{H}_{v, 3})} (|h| - 1) = \sum_{u \in N(v)} |\{h \in E(\mathcal{H}_{v, 3}): u \in h\}| \leq (t-1) d(v).
            \end{equation}
            % For each hyperedge $h$ in $\mathcal{H}_{v, 1} \cup \mathcal{H}_{v, 2}$, the size of $h \cap N(v)$ is at most $t$.
            % Let $S \subseteq N(v)$ be a set of size $t$.

            % \begin{claimthm}
            %     If $S$ is contained in at least $t$ hyperedges of $\mathcal{H}_{v, 2}$, then no vertex in $S$ belongs to any other hyperedge $g$, $|g \cap S| \leq t-1$, of $\mathcal{H}_{v, 2}$ besides these $t$ hyperedges.
            % \end{claimthm}
            % \begin{proof}
            %     Suppose not, let $h_1, h_2, \dots, h_t$ be $t$ hyperedges containing $S = \{u_1, \dots, u_t\}$.
            %     Assume that $u_1$ is also contained in $h_{t+1}$, $|h_{t+1} \cap S| \leq t-1$.
            %     Let $w_i = h_i \setminus S$, $u_i \in h_i$ such that $u_i \neq u_j$ for $i \neq j$.
            %     Then there are $t$ copies of $C_4$ $v u_1 w_2 u_2$, $v u_1 w_3 u_3$, $\dots$, $v u_1 w_t u_t$ and $v u_1 w_{t+1} u_{t+1}$ which form a $B_t$, a contradiction.
            % \end{proof}

            % \begin{claimthm}
            %     If $S$ is contained in at least $1$ and at most $t-1$ hyperedges of $\mathcal{H}_{v, 2}$, then each vertex in $S$ is contained in at most $t-1$ hyperedges of $\mathcal{H}_{v, 2}$.
            % \end{claimthm}
            % \begin{proof}
            %     Suppose not, let $h_1, h_2, \dots, h_{t_1}$ be $t_1$ hyperedges containing $S = \{u_1, u_2, \dots, u_t\}$, where $1 \leq t_1 \leq t-1$.
            %     Assume that $u_1$ is also contained in $h_{t_1 + 1}, \dots, h_t$.
            %     Let $w_i = h_i \setminus S$, $u_i \in h_i$ such that $u_i \neq u_j$ for $i \neq j$.
            %     Then there are $t$ copies of $C_4$ $v u_1 w_2 u_2$, $v u_1 w_3 u_3$, $\dots$, $v u_1 w_t u_t$ and $v u_1 w_{t+1} u_{t+1}$ which form a $B_t$, a contradiction.
            % \end{proof}
            % This means $\sum_{h \in E(\mathcal{H}_{v, 2}), d(h \cap N(v)) \leq t-1} (|h| - 1) \leq (t-1)d(v).$
            % The $d(S) - 1$ is the number of common neighbors of $S$ besides $v$.
            % Suppose that there are $x \leq n_1$ hyperedges $h$ of $\mathcal{H}_{v, 2}$ such that $d_{\mathcal{H}_{v, 2}}(h \cap N(v)) \geq t$.
            % Thus $\sum_{h \in E(\mathcal{H}_{v, 2}), d(h \cap N(v)) \geq t} (|h| - 1)  \leq  \sum_{\substack{S \subseteq N(v), |S| = t, S \text{ is good}}} (d(S) - 1) + (t-1)x$.
            % There are at most $n_1-x$ hypereges in $\mathcal{H}_{v, 1}$, so $\sum_{h \in E(\mathcal{H}_{v, 1})} (|h| - 1) \leq (t-1)(n_1 - x)$.
            Notice that $v \in V_1$,
            by the construction of $\mathcal{H}_v$, each hyperedge in $\mathcal{H}_v$ consists of a vertex in $V_1$ together with its neighbors.
            Thus, there are at most $n_1$ hyperedges in $\mathcal{H}_v$.
            Suppose that there are $x \leq n_1$ hyperedges $h$ of $\mathcal{H}_{v, 2}$ in $E(\mathcal{H}_{v, 2}) \setminus E_2$.
            Thus $\sum_{h \in E(\mathcal{H}_{v, 2}) \setminus E_2} (|h| - 1)  \leq  \sum_{\substack{S \subseteq N(v), |S| = t, S \text{ is good}}} (d(S) - 1) + (t-1)x$.
            There are at most $n_1-x$ hyperedges in $\mathcal{H}_{v, 1}$, so $\sum_{h \in E(\mathcal{H}_{v, 1})} (|h| - 1) \leq (t-1)(n_1 - x)$.
            As in the proof of Lemma \ref{lem: lower bound of number of C4 for Bt}, we have
            \begin{equation}\label{eq: lem2.2 2}
                \sum_{i = 1}^{2} \sum_{h \in E(\mathcal{H}_{v, i})} (|h| - 1) \leq (t-1)d(v) + (t-1)n_1 + \sum_{\substack{S \subseteq N(v), |S| = t, \\ S \text{ is good}}} (d(S) - 1).
            \end{equation}

            Combining Equation (\ref{eq: lem2.2 1}), (\ref{eq: lem2.2 2}) and (\ref{eq: sum of dv bip}), we have
            \begin{equation*}
                \begin{aligned}
                \sum_{\substack{S \subseteq N(v), |S| = t \\ S \text{ is good}}} (d(S) - 1)
                \geq& \sum_{u \in N(v)} d(u) - d(v) - (t-1) d(v) - (t-1) d(v) - (t-1) n_1 \\
                = & \sum_{u \in N(v)} d(u) - (2t-1) d(v) - (t-1) n_1.
                \end{aligned}
            \end{equation*}
            Summing over all $v \in V_1$, we have
            \begin{equation*}
                \begin{aligned}
                    & \sum_{v \in V_1} \sum_{\substack{S \subseteq N(v), |S| = t \\ S \text{ is good}}} (d(S) - 1) \\
                    \geq& \sum_{v \in V_1} \left(\sum_{u \in N(v)} d(u) - (2t-1) d(v) - (t-1) n_1 \right) \\
                    =& \sum_{u \in V_2} d(u)^2 - (2t-1) e(G) - (t-1) n_1^2.
                \end{aligned}
            \end{equation*}
            On the other hand, each set $S$ of size $t$ can be contained in $d(S)$ neighborhoods. So
            \begin{equation*}
                \sum_{\substack{S \subseteq V_2, |S| = t \\ S \text{ is good}}} d(S)(d(S) - 1) = \sum_{v \in V_1} \sum_{\substack{S \subseteq N(v), |S| = t \\ S \text{ is good}}} (d(S) - 1).
            \end{equation*}
            Thus, we have
            \begin{equation}\label{eq: lower bound of number of C4 1}
                \sum_{\substack{S \subseteq V_2, |S| = t \\ S \text{ is good}}} d(S)(d(S) - 1) \geq \sum_{u \in V_2} d(u)^2 - (2t-1) e(G) - (t-1) n_1^2.
            \end{equation}

            For each set $S \subseteq V_2$ of size $t$, any two common neighbors of $S$ together with two vertices in $S$ form a $C_4$.
            Thus
            \begin{equation}\label{eq: sum of dv 2}
                N_2 (C_4, G) = \sum_{\substack{S \subseteq V_2, |S| = t \\ S \text{ is good}}} \binom{d(S)}{2} \binom{t}{2}.
            \end{equation}

            Combining Equation (\ref{eq: lower bound of number of C4 1}) and Equation (\ref{eq: sum of dv 2}), we have
            \begin{equation*}
                N_2(C_4, G) \geq \frac{1}{2} \binom{t}{2} \left(\sum_{u \in V_2} d(u)^2 - (2t-1) e(G) - (t-1) n_1^2 \right).
            \end{equation*}

            Similarly to the above, we obtain $N_1(C_4, G) \geq \frac{1}{2} \binom{t}{2} (\sum_{u \in V_1} d(u)^2 - (2t-1) e(G) - (t-1) n_2^2)$.
            The proof is complete.
        \end{proof}

\end{document}
